\IfFileExists{stacks-project.cls}{%
\documentclass{stacks-project}
}{%
\documentclass{amsart}
}

% For dealing with references we use the comment environment
\usepackage{verbatim}
\newenvironment{reference}{\comment}{\endcomment}
%\newenvironment{reference}{}{}
\newenvironment{slogan}{\comment}{\endcomment}
\newenvironment{history}{\comment}{\endcomment}

% For commutative diagrams we use Xy-pic
\usepackage[all]{xy}

% We use 2cell for 2-commutative diagrams.
\xyoption{2cell}
\UseAllTwocells

% We use multicol for the list of chapters between chapters
\usepackage{multicol}

% This is generally recommended for better output
\usepackage{lmodern}
\usepackage[T1]{fontenc}

% For cross-file-references
\usepackage{xr-hyper}

% Package for hypertext links:
\usepackage{hyperref}

% For any local file, say "hello.tex" you want to link to please
% use \externaldocument[hello-]{hello}
\externaldocument[introduction-]{introduction}
\externaldocument[conventions-]{conventions}
\externaldocument[sets-]{sets}
\externaldocument[categories-]{categories}
\externaldocument[topology-]{topology}
\externaldocument[sheaves-]{sheaves}
\externaldocument[sites-]{sites}
\externaldocument[stacks-]{stacks}
\externaldocument[fields-]{fields}
\externaldocument[algebra-]{algebra}
\externaldocument[brauer-]{brauer}
\externaldocument[homology-]{homology}
\externaldocument[derived-]{derived}
\externaldocument[simplicial-]{simplicial}
\externaldocument[more-algebra-]{more-algebra}
\externaldocument[smoothing-]{smoothing}
\externaldocument[modules-]{modules}
\externaldocument[sites-modules-]{sites-modules}
\externaldocument[injectives-]{injectives}
\externaldocument[cohomology-]{cohomology}
\externaldocument[sites-cohomology-]{sites-cohomology}
\externaldocument[dga-]{dga}
\externaldocument[dpa-]{dpa}
\externaldocument[sdga-]{sdga}
\externaldocument[hypercovering-]{hypercovering}
\externaldocument[schemes-]{schemes}
\externaldocument[constructions-]{constructions}
\externaldocument[properties-]{properties}
\externaldocument[morphisms-]{morphisms}
\externaldocument[coherent-]{coherent}
\externaldocument[divisors-]{divisors}
\externaldocument[limits-]{limits}
\externaldocument[varieties-]{varieties}
\externaldocument[topologies-]{topologies}
\externaldocument[descent-]{descent}
\externaldocument[perfect-]{perfect}
\externaldocument[more-morphisms-]{more-morphisms}
\externaldocument[flat-]{flat}
\externaldocument[groupoids-]{groupoids}
\externaldocument[more-groupoids-]{more-groupoids}
\externaldocument[etale-]{etale}
\externaldocument[chow-]{chow}
\externaldocument[intersection-]{intersection}
\externaldocument[pic-]{pic}
\externaldocument[weil-]{weil}
\externaldocument[adequate-]{adequate}
\externaldocument[dualizing-]{dualizing}
\externaldocument[duality-]{duality}
\externaldocument[discriminant-]{discriminant}
\externaldocument[derham-]{derham}
\externaldocument[local-cohomology-]{local-cohomology}
\externaldocument[algebraization-]{algebraization}
\externaldocument[curves-]{curves}
\externaldocument[resolve-]{resolve}
\externaldocument[models-]{models}
\externaldocument[functors-]{functors}
\externaldocument[equiv-]{equiv}
\externaldocument[pione-]{pione}
\externaldocument[etale-cohomology-]{etale-cohomology}
\externaldocument[proetale-]{proetale}
\externaldocument[relative-cycles-]{relative-cycles}
\externaldocument[more-etale-]{more-etale}
\externaldocument[trace-]{trace}
\externaldocument[crystalline-]{crystalline}
\externaldocument[spaces-]{spaces}
\externaldocument[spaces-properties-]{spaces-properties}
\externaldocument[spaces-morphisms-]{spaces-morphisms}
\externaldocument[decent-spaces-]{decent-spaces}
\externaldocument[spaces-cohomology-]{spaces-cohomology}
\externaldocument[spaces-limits-]{spaces-limits}
\externaldocument[spaces-divisors-]{spaces-divisors}
\externaldocument[spaces-over-fields-]{spaces-over-fields}
\externaldocument[spaces-topologies-]{spaces-topologies}
\externaldocument[spaces-descent-]{spaces-descent}
\externaldocument[spaces-perfect-]{spaces-perfect}
\externaldocument[spaces-more-morphisms-]{spaces-more-morphisms}
\externaldocument[spaces-flat-]{spaces-flat}
\externaldocument[spaces-groupoids-]{spaces-groupoids}
\externaldocument[spaces-more-groupoids-]{spaces-more-groupoids}
\externaldocument[bootstrap-]{bootstrap}
\externaldocument[spaces-pushouts-]{spaces-pushouts}
\externaldocument[spaces-chow-]{spaces-chow}
\externaldocument[groupoids-quotients-]{groupoids-quotients}
\externaldocument[spaces-more-cohomology-]{spaces-more-cohomology}
\externaldocument[spaces-simplicial-]{spaces-simplicial}
\externaldocument[spaces-duality-]{spaces-duality}
\externaldocument[formal-spaces-]{formal-spaces}
\externaldocument[restricted-]{restricted}
\externaldocument[spaces-resolve-]{spaces-resolve}
\externaldocument[formal-defos-]{formal-defos}
\externaldocument[defos-]{defos}
\externaldocument[cotangent-]{cotangent}
\externaldocument[examples-defos-]{examples-defos}
\externaldocument[algebraic-]{algebraic}
\externaldocument[examples-stacks-]{examples-stacks}
\externaldocument[stacks-sheaves-]{stacks-sheaves}
\externaldocument[criteria-]{criteria}
\externaldocument[artin-]{artin}
\externaldocument[quot-]{quot}
\externaldocument[stacks-properties-]{stacks-properties}
\externaldocument[stacks-morphisms-]{stacks-morphisms}
\externaldocument[stacks-limits-]{stacks-limits}
\externaldocument[stacks-cohomology-]{stacks-cohomology}
\externaldocument[stacks-perfect-]{stacks-perfect}
\externaldocument[stacks-introduction-]{stacks-introduction}
\externaldocument[stacks-more-morphisms-]{stacks-more-morphisms}
\externaldocument[stacks-geometry-]{stacks-geometry}
\externaldocument[moduli-]{moduli}
\externaldocument[moduli-curves-]{moduli-curves}
\externaldocument[examples-]{examples}
\externaldocument[exercises-]{exercises}
\externaldocument[guide-]{guide}
\externaldocument[desirables-]{desirables}
\externaldocument[coding-]{coding}
\externaldocument[obsolete-]{obsolete}
\externaldocument[fdl-]{fdl}
\externaldocument[index-]{index}

% Theorem environments.
%
\theoremstyle{plain}
\newtheorem{theorem}[subsection]{Theorem}
\newtheorem{proposition}[subsection]{Proposition}
\newtheorem{lemma}[subsection]{Lemma}

\theoremstyle{definition}
\newtheorem{definition}[subsection]{Definition}
\newtheorem{example}[subsection]{Example}
\newtheorem{exercise}[subsection]{Exercise}
\newtheorem{situation}[subsection]{Situation}

\theoremstyle{remark}
\newtheorem{remark}[subsection]{Remark}
\newtheorem{remarks}[subsection]{Remarks}

\numberwithin{equation}{subsection}

% Macros
%
\def\lim{\mathop{\mathrm{lim}}\nolimits}
\def\colim{\mathop{\mathrm{colim}}\nolimits}
\def\Spec{\mathop{\mathrm{Spec}}}
\def\Hom{\mathop{\mathrm{Hom}}\nolimits}
\def\Ext{\mathop{\mathrm{Ext}}\nolimits}
\def\SheafHom{\mathop{\mathcal{H}\!\mathit{om}}\nolimits}
\def\SheafExt{\mathop{\mathcal{E}\!\mathit{xt}}\nolimits}
\def\Sch{\mathit{Sch}}
\def\Mor{\mathop{\mathrm{Mor}}\nolimits}
\def\Ob{\mathop{\mathrm{Ob}}\nolimits}
\def\Sh{\mathop{\mathit{Sh}}\nolimits}
\def\NL{\mathop{N\!L}\nolimits}
\def\CH{\mathop{\mathrm{CH}}\nolimits}
\def\proetale{{pro\text{-}\acute{e}tale}}
\def\etale{{\acute{e}tale}}
\def\QCoh{\mathit{QCoh}}
\def\Ker{\mathop{\mathrm{Ker}}}
\def\Im{\mathop{\mathrm{Im}}}
\def\Coker{\mathop{\mathrm{Coker}}}
\def\Coim{\mathop{\mathrm{Coim}}}

% Boxtimes
%
\DeclareMathSymbol{\boxtimes}{\mathbin}{AMSa}{"02}

%
% Macros for moduli stacks/spaces
%
\def\QCohstack{\mathcal{QC}\!\mathit{oh}}
\def\Cohstack{\mathcal{C}\!\mathit{oh}}
\def\Spacesstack{\mathcal{S}\!\mathit{paces}}
\def\Quotfunctor{\mathrm{Quot}}
\def\Hilbfunctor{\mathrm{Hilb}}
\def\Curvesstack{\mathcal{C}\!\mathit{urves}}
\def\Polarizedstack{\mathcal{P}\!\mathit{olarized}}
\def\Complexesstack{\mathcal{C}\!\mathit{omplexes}}
% \Pic is the operator that assigns to X its picard group, usage \Pic(X)
% \Picardstack_{X/B} denotes the Picard stack of X over B
% \Picardfunctor_{X/B} denotes the Picard functor of X over B
\def\Pic{\mathop{\mathrm{Pic}}\nolimits}
\def\Picardstack{\mathcal{P}\!\mathit{ic}}
\def\Picardfunctor{\mathrm{Pic}}
\def\Deformationcategory{\mathcal{D}\!\mathit{ef}}
\usepackage{cleveref}
\crefname{originallemma}{original lemma}{original lemmas}
\Crefname{originallemma}{Original lemma}{Original lemmas}
\makeatletter
\newlabel{lemma-group-over-field-locally-finite-type-dimension}{{06FE}{source}{}{sourcepointer-1}{}}
\newlabel{lemma-group-over-field-locally-finite-type-dimension@cref}{{[originallemma][0][]06FE}{[1][0][]source}{}{}{}}
\makeatother
\crefname{originallemma}{original lemma}{original lemmas}
\Crefname{originallemma}{Original lemma}{Original lemmas}
\makeatletter
\newlabel{lemma-groupoid-on-field-locally-finite-type-dimension}{{04LU}{source}{}{sourcepointer-2}{}}
\newlabel{lemma-groupoid-on-field-locally-finite-type-dimension@cref}{{[originallemma][0][]04LU}{[1][0][]source}{}{}{}}
\makeatother
\crefname{originalreference}{original reference}{original references}
\Crefname{originalreference}{Original reference}{Original references}
\makeatletter
\newlabel{morphisms-lemma-finite-type-noetherian}{{01T6}{source}{}{sourcepointer-3}{}}
\newlabel{morphisms-lemma-finite-type-noetherian@cref}{{[originalreference][0][]01T6}{[1][0][]source}{}{}{}}
\makeatother
\crefname{originalreference}{original reference}{original references}
\Crefname{originalreference}{Original reference}{Original references}
\makeatletter
\newlabel{properties-lemma-Noetherian-topology}{{01OZ}{source}{}{sourcepointer-4}{}}
\newlabel{properties-lemma-Noetherian-topology@cref}{{[originalreference][0][]01OZ}{[1][0][]source}{}{}{}}
\makeatother
\crefname{originalreference}{original reference}{original references}
\Crefname{originalreference}{Original reference}{Original references}
\makeatletter
\newlabel{situation-groupoid-on-field}{{06DY}{source}{}{sourcepointer-5}{}}
\newlabel{situation-groupoid-on-field@cref}{{[originalreference][0][]06DY}{[1][0][]source}{}{}{}}
\makeatother
\crefname{originalreference}{original reference}{original references}
\Crefname{originalreference}{Original reference}{Original references}
\makeatletter
\newlabel{spaces-morphisms-lemma-composition-finite-type}{{03XG}{source}{}{sourcepointer-6}{}}
\newlabel{spaces-morphisms-lemma-composition-finite-type@cref}{{[originalreference][0][]03XG}{[1][0][]source}{}{}{}}
\makeatother
\crefname{originalreference}{original reference}{original references}
\Crefname{originalreference}{Original reference}{Original references}
\makeatletter
\newlabel{spaces-morphisms-lemma-etale-locally-finite-presentation}{{0468}{source}{}{sourcepointer-7}{}}
\newlabel{spaces-morphisms-lemma-etale-locally-finite-presentation@cref}{{[originalreference][0][]0468}{[1][0][]source}{}{}{}}
\makeatother
\crefname{originalreference}{original reference}{original references}
\Crefname{originalreference}{Original reference}{Original references}
\makeatletter
\newlabel{spaces-morphisms-lemma-finite-presentation-finite-type}{{0464}{source}{}{sourcepointer-8}{}}
\newlabel{spaces-morphisms-lemma-finite-presentation-finite-type@cref}{{[originalreference][0][]0464}{[1][0][]source}{}{}{}}
\makeatother
\crefname{originalreference}{original reference}{original references}
\Crefname{originalreference}{Original reference}{Original references}
\makeatletter
\newlabel{topology-lemma-locally-Noetherian-locally-connected}{{04MF}{source}{}{sourcepointer-9}{}}
\newlabel{topology-lemma-locally-Noetherian-locally-connected@cref}{{[originalreference][0][]04MF}{[1][0][]source}{}{}{}}
\makeatother
\crefname{originalreference}{original reference}{original references}
\Crefname{originalreference}{Original reference}{Original references}
\makeatletter
\newlabel{varieties-proposition-unique-base-field}{{04MK}{source}{}{sourcepointer-10}{}}
\newlabel{varieties-proposition-unique-base-field@cref}{{[originalreference][0][]04MK}{[1][0][]source}{}{}{}}
\makeatother
\title{A groupoid over a field has the dimension of its stabilizer}
\author{The Stacks Project Authors}
\date{stacks-project:spaces-more-groupoids.tex:local-snapshot-20261012; GNU FDL 1.2 or later}
\begin{document}
\maketitle
\noindent\textit{Editorial note.} No omission markers in-span. The equidimensionality input from the uncarded group-side dimension lemma is declared as external.\par\medskip
\noindent\textit{Allowed standard prerequisites (editorial).}\begin{itemize}
\item Equidimensionality of R and G. lemma-groupoid-on-field-locally-finite-type-dimension (same chapter, carded L26644) and lemma-group-over-field-locally-finite-type-dimension (same chapter, NOT carded, declared external). Reduces the claim to dim\_e(R) = dim\_e(G).
\item Unique base field in a variety ring. Varieties, proposition unique-base-field (cross-chapter, named). Identifies the integral closures k\_a and k\_b of the two field images.
\end{itemize}\medskip
\ifdefined\hypertarget\hypertarget{sourcepointer-1}{}\fi
\noindent\textit{Source reference.} Original lemma 06FE: sources/stacks-spaces-more-groupoids/pinned-original-source.tex, source lines 1160--1160. Navigation pointer only: the counted proof cites this statement by \textbackslash{}ref as a standing input; the deduction the counted proof has to make is carried out inside the counted span, and no part of that deduction is delegated to this pointer.\par\medskip
\ifdefined\hypertarget\hypertarget{sourcepointer-2}{}\fi
\noindent\textit{Source reference.} Original lemma 04LU: sources/stacks-spaces-more-groupoids/pinned-original-source.tex, source lines 1107--1107. Navigation pointer only: the counted proof cites this statement by \textbackslash{}ref as a standing input; the deduction the counted proof has to make is carried out inside the counted span, and no part of that deduction is delegated to this pointer.\par\medskip
\ifdefined\hypertarget\hypertarget{sourcepointer-3}{}\fi
\noindent\textit{Source reference.} Original reference 01T6: focused-seeds/geometry/S9-smg-groupoid-on-field-dimension-equal-stabilizer/ref-tags.txt, source lines 5--5. Navigation pointer only: the counted proof cites this statement by \textbackslash{}ref as a standing input; the deduction the counted proof has to make is carried out inside the counted span, and no part of that deduction is delegated to this pointer.\par\medskip
\ifdefined\hypertarget\hypertarget{sourcepointer-4}{}\fi
\noindent\textit{Source reference.} Original reference 01OZ: focused-seeds/geometry/S9-smg-groupoid-on-field-dimension-equal-stabilizer/ref-tags.txt, source lines 6--6. Navigation pointer only: the counted proof cites this statement by \textbackslash{}ref as a standing input; the deduction the counted proof has to make is carried out inside the counted span, and no part of that deduction is delegated to this pointer.\par\medskip
\ifdefined\hypertarget\hypertarget{sourcepointer-5}{}\fi
\noindent\textit{Source reference.} Original reference 06DY: sources/stacks-spaces-more-groupoids/pinned-original-source.tex, source lines 859--859. Navigation pointer only: the counted proof cites this statement by \textbackslash{}ref as a standing input; the deduction the counted proof has to make is carried out inside the counted span, and no part of that deduction is delegated to this pointer.\par\medskip
\ifdefined\hypertarget\hypertarget{sourcepointer-6}{}\fi
\noindent\textit{Source reference.} Original reference 03XG: focused-seeds/geometry/S9-smg-groupoid-on-field-dimension-equal-stabilizer/ref-tags.txt, source lines 7--7. Navigation pointer only: the counted proof cites this statement by \textbackslash{}ref as a standing input; the deduction the counted proof has to make is carried out inside the counted span, and no part of that deduction is delegated to this pointer.\par\medskip
\ifdefined\hypertarget\hypertarget{sourcepointer-7}{}\fi
\noindent\textit{Source reference.} Original reference 0468: focused-seeds/geometry/S9-smg-groupoid-on-field-dimension-equal-stabilizer/ref-tags.txt, source lines 8--8. Navigation pointer only: the counted proof cites this statement by \textbackslash{}ref as a standing input; the deduction the counted proof has to make is carried out inside the counted span, and no part of that deduction is delegated to this pointer.\par\medskip
\ifdefined\hypertarget\hypertarget{sourcepointer-8}{}\fi
\noindent\textit{Source reference.} Original reference 0464: focused-seeds/geometry/S9-smg-groupoid-on-field-dimension-equal-stabilizer/ref-tags.txt, source lines 9--9. Navigation pointer only: the counted proof cites this statement by \textbackslash{}ref as a standing input; the deduction the counted proof has to make is carried out inside the counted span, and no part of that deduction is delegated to this pointer.\par\medskip
\ifdefined\hypertarget\hypertarget{sourcepointer-9}{}\fi
\noindent\textit{Source reference.} Original reference 04MF: focused-seeds/geometry/S9-smg-groupoid-on-field-dimension-equal-stabilizer/ref-tags.txt, source lines 10--10. Navigation pointer only: the counted proof cites this statement by \textbackslash{}ref as a standing input; the deduction the counted proof has to make is carried out inside the counted span, and no part of that deduction is delegated to this pointer.\par\medskip
\ifdefined\hypertarget\hypertarget{sourcepointer-10}{}\fi
\noindent\textit{Source reference.} Original reference 04MK: focused-seeds/geometry/S9-smg-groupoid-on-field-dimension-equal-stabilizer/ref-tags.txt, source lines 11--11. Navigation pointer only: the counted proof cites this statement by \textbackslash{}ref as a standing input; the deduction the counted proof has to make is carried out inside the counted span, and no part of that deduction is delegated to this pointer.\par\medskip
\begin{lemma}
\label{lemma-groupoid-on-field-dimension-equal-stabilizer}
In
Situation \ref{situation-groupoid-on-field}
assume $s, t$ are locally of finite type.
Let
$G = \Spec(k)
\times_{\Delta, \Spec(k) \times_B \Spec(k), t \times s} R$
be the stabilizer group algebraic space.
Then we have $\dim(R) = \dim(G)$.
\end{lemma}

\begin{proof}
Since $G$ and $R$ are equidimensional (see
Lemmas \ref{lemma-groupoid-on-field-locally-finite-type-dimension} and
\ref{lemma-group-over-field-locally-finite-type-dimension})
it suffices to prove that $\dim_e(R) = \dim_e(G)$. Let $V$ be an affine scheme,
$v \in V$, and let $\varphi : V \to R$ be an \'etale morphism of schemes
such that $\varphi(v) = e$. Note that $V$ is a Noetherian scheme as
$s \circ \varphi$ is locally of finite type as a composition of morphisms
locally of finite type and as $V$ is quasi-compact (use
Morphisms of Spaces, Lemmas
\ref{spaces-morphisms-lemma-composition-finite-type},
\ref{spaces-morphisms-lemma-etale-locally-finite-presentation}, and
\ref{spaces-morphisms-lemma-finite-presentation-finite-type}
and
Morphisms, Lemma \ref{morphisms-lemma-finite-type-noetherian}).
Hence $V$ is locally connected (see
Properties, Lemma \ref{properties-lemma-Noetherian-topology}
and
Topology, Lemma \ref{topology-lemma-locally-Noetherian-locally-connected}).
Thus we may replace $V$ by the connected component containing $v$ (it
is still affine as it is an open and closed subscheme of $V$).
Set $T = V_{red}$ equal to the reduction of $V$. Consider the two
morphisms $a, b : T \to \Spec(k)$ given by
$a = s \circ \varphi|_T$ and $b = t \circ \varphi|_T$. Note that
$a, b$ induce the same field map $k \to \kappa(v)$ because $\varphi(v) = e$!
Let $k_a \subset \Gamma(T, \mathcal{O}_T)$ be the integral closure of
$a^\sharp(k) \subset \Gamma(T, \mathcal{O}_T)$. Similarly, let
$k_b \subset \Gamma(T, \mathcal{O}_T)$ be the integral closure of
$b^\sharp(k) \subset \Gamma(T, \mathcal{O}_T)$. By
Varieties, Proposition \ref{varieties-proposition-unique-base-field}
we see that $k_a = k_b$. Thus we obtain the following commutative diagram
$$
\xymatrix{
k \ar[rd]^a \ar[rrrd] \\
& k_a = k_b \ar[r] & \Gamma(T, \mathcal{O}_T) \ar[r] & \kappa(v) \\
k \ar[ru]_b \ar[rrru]
}
$$
As discussed above the long arrows are equal.
Since $k_a = k_b \to \kappa(v)$ is injective we conclude that
the two morphisms $a$ and $b$ agree. Hence $T \to R$ factors through $G$.
It follows that $R_{red} = G_{red}$ in an open neighbourhood of $e$
which certainly implies that $\dim_e(R) = \dim_e(G)$.
\end{proof}

\end{document}
